\chapter{Routines for interfacing with other codes: \\EIRCOP}
\label{sec4}

\medskip

\noindent

{\bf General remarks}

\medskip  %

\noindent

The ``background'' or ``host'' medium, usually the ``plasma'' is fixed in an EIRENE run, and usually either specified
in input block 5 (\ref{sec2.5}), or the background volume tallies are simply ready from data files.

However the mutual coupling between test-particle species treated by EIRENE and the background medium may be very
strong, even if the test particle density is quite low compared to the background (plasma) density. This is
because in-elastic collisions may lead to strong particle, momentum and/or energy sources for the balances governing the
background medium, and these ``chemical sources" may even be the dominant terms for the global flow problem,
at least in some parts of the computational volume (e.g.
typically in the divertor of fusion reactors). Original diffusion-advection equations for a background medium
there may turn into diffusion-advection-reaction type equations and even change their mathematical characteristics.

For such type of ``background medium models'' (often ``Navier Stokes like" or related) then special interfacing procedures
between the kinetic Monte Carlo component (EIRENE) and the numerical background model are required.

For the EIRENE code such procedures have been developed in the late eighties, originally for the SOLXY 2D plasma code (\cite{kn:Gerhauser88a})
but have then been transferred and extended for the rather well established (at that time) 2D multi-fluid plasma solver B2.

The EIRENE modules developed for code interfacing are described in this section, original references are: Reiter, \cite{kn:Net}, and first applications \cite{kn:Reiter91b} (1991)
and the most detailed discussion probably still is the FZ J\"ulich report JUEL-2872,
Maddison, Reiter, \cite{kn:Maddison94},  (1994).

Special Monte Carlo options, such as correlated sampling 
(see flag NLCRR in input block 1, \ref{sec2.1}) or the semi-implicit
coupling schemes based on source term adaptations/rescalings 
to reduce Monte Carlo run-time in these iterative non-linear schemes
(jargon: ``short cycles"), have also had their
origin in those early studies, loc cit., but have meanwhile been generalized significantly, e.g. to allow selective
``short cycling" for individual strata,
while for others (notably for volume recombination sources)
full new Monte Carlo estimates are retained.




The names of all routines in the interfacing block
end with ...COP.

\section{Routine for interfacing INFCOP} \label{sec4.1}


To write an interfacing routine INFCOP
in order to couple EIRENE to
another code is already a quite formidable task, a new scientific challenge and it is highly
recommended that, in addition to the information given below,
close interaction with scientists at FZ Juelich may be needed. One also might
ask for a few sample routines INFCOP, which are available from
the authors.

In the rest of this section plasma data and geometrical (grid) information, from interfacing with an external model/code are
described with reference to their location and size on a single large, 1D storage array \texttt{RWK}.
This array \texttt{RWK} which is then used in the initialization phase of an EIRENE run to bring this information into
regular EIRENE storage places/names.

In Version from 2001 or later, rather than filling 1D array \texttt{RWK}, now directly
input tallies TEINTF, TIINTF, DIINTF, ... etc. can be filled. These are pointers now to a vector (now named
\texttt{PLASMA\_BACKGROUND}), but which has the same structure as the old \texttt{RWK} array. Apart from that the rest remains the same
and the material in this section is still valid.

Data are transfered from subroutine INFCOP into EIRENE via the
EIRENE work array \texttt{RWK}  (Module CSPEI).
They are read onto the EIRENE arrays for input volume tallies (\ref{tab1}) by calls of subroutine PROFR in
the initialization phase individually for each tally,
for which the flag INDPRO has the value 6 or 7 for a particular input tally
(see section \ref{sec2.5}).
The following addresses are foreseen on \texttt{RWK} for this data transferring
procedure:


{\sl Plasma data (INDPRO = 6 option) }

\begin{verbatim}
TEIN  :  (RWK ((0         ) * NRAD + J), J=1,NSBOX)
TIIN  :  (RWK ((1 + 0*NPLS) * NRAD + J), J=1,NSBOX,I=1,NPLS)
DIIN  :  (RWK ((1 + 1*NPLS) * NRAD + J), J=1,NSBOX,I=1,NPLS)
VXIN  :  (RWK ((1 + 2*NPLS) * NRAD + J), J=1,NSBOX,I=1,NPLS)
VYIN  :  (RWK ((1 + 3*NPLS) * NRAD + J), J=1,NSBOX,I=1,NPLS)
VZIN  :  (RWK ((1 + 4*NPLS) * NRAD + J), J=1,NSBOX,I=1,NPLS)
BXIN  :  (RWK ((1 + 5*NPLS) * NRAD + J), J=1,NSBOX)
BYIN  :  (RWK ((2 + 5*NPLS) * NRAD + J), J=1,NSBOX)
BZIN  :  (RWK ((3 + 5*NPLS) * NRAD + J), J=1,NSBOX)
BFIN  :  (RWK ((4 + 5*NPLS) * NRAD + J), J=1,NSBOX)
VLIN  :  (RWK ((5 + 5*NPLS) * NRAD + J), J=1,NSBOX)
ADIN  :  (RWK ((6 + 5*NPLS) * NRAD + J), J=1,NSBOX,I=1,NAIN)
\end{verbatim}
{\sl Geometrical data (INDGRD = 6 option) }
\begin{verbatim}
LEVGEO = 1 or LEVGEO = 2:
RSURF :  (RWK ((6+5*NPLS+NAIN) * NRAD +                 J), J=1,N1ST)
EP    :  (RWK ((6+5*NPLS+NAIN) * NRAD +   N1ST +        J), J=1,N1ST)
EL    :  (RWK ((6+5*NPLS+NAIN) * NRAD + 2*N1ST +        J), J=1,N1ST)
TR    :  (RWK ((6+5*NPLS+NAIN) * NRAD + 3*N1ST +        J), J=1,N1ST)
PSURF :  (RWK ((6+5*NPLS+NAIN) * NRAD + 4*N1ST +        J), J=1,N2ND)
TSURF :  (RWK ((6+5*NPLS+NAIN) * NRAD + 4*N1ST + N2ND + J), J=1,N3RD)
LEVGEO = 3
PGINTF:  (RWK ((6+5*NPLS+NAIN) * NRAD + J), J=1,NPMAX)
TSURF :  (RWK ((6+5*NPLS+NAIN) * NRAD + NPMAX + J), J=1,N3RD)
LEVGEO = 4
TRINTF:  (RWK ((6+5*NPLS+NAIN) * NRAD + J), J=1,NTMAX)
TSURF :  (RWK ((6+5*NPLS+NAIN) * NRAD + NTMAX + J), J=1,N3RD)
LEVGEO = 5
THINTF:  (RWK ((6+5*NPLS+NAIN) * NRAD + J), J=1,NHMAX)
\end{verbatim}

\subsection{entry IF0COP} \label{sec4.1.0}
{\sl Geometry data (INDGRD = 6 option)}

\medskip

In the $LEVGEO=3$ option (2D grid of quadrangles) PGINTF is an array of length NPMAX,
which contains all relevant
information to generate a 2D mesh of polygons in the x-y plane.
It is equivalenced to the common block CPOLYG
via the statement:

EQUIVALENCE (PGINTF(1),XPLG(1,1)),...

In the $LEVGEO=4$ option (2D grid of triangles) TRINTF is an array of length NTMAX,
which contains all relevant
information to generate a mesh of triangles in the x-y plane.
It is equivalenced to the common block CTRIA
via the statement:

EQUIVALENCE (TRINTF(1),XT(1)),...

In the $LEVGEO=5$ option (3D grid of tetrahedons) THINTF is an array of length NHMAX, which
contains all relevant information to generate a mesh of
tetrahedrons in the 3d computational domain. It is equivalenced to
the common block CTETRA via the statement:

EQUIVALENCE (THINTF(1),XT(1)),...

\subsection{entry IF1COP} \label{sec4.1.1}
{\sl Plasma (background) data (INDPRO = 6 option)}


\medskip


to be written
\subsection{entry IF2COP(ISTRA)} \label{sec4.1.2}
{\sl Primary source data (INDSRC = 6 option)}

\medskip
In case of an EIRENE run, in which input information is obtained
from an external code (e.g.: B2, EMC3, DIVIMP, ...), the
stratification of the primary source should be such that the first
NTARGI (see \ref{sec2.14}) strata are determined by the plasma
fluxes onto surfaces (``recycling surface sources"). In this case
the surface source distribution $Q_S(\vr,\vv,i,t)$ (section \ref{sec1.2}) can be defined automatically from the interfacing
routine, using the data specified in input block 14. The remaining
primary sources NTARGI+1, ..., NSTRAI (e.g. gas puff, volume
recombination sources, etc...) still have to defined in input
block 7.  Whether input block 7 or input block 14 is used for
a particular surface source ISTRA $\le$ NTARGI is controlled by
the flag INDSRC.  For sources ISTRA $>$ NTARGI the input flag INDSRC  is irrelevant.

if INDSRC(ISTRA) $<$ 0 then interfacing routine IF2COP is not called.
Input data are used from block 7, see \ref{sec2.7}.

\medskip

if INDSRC(ISTRA) = 0--5, then the input flags described in section
\ref{sec2.7} are used to define the spatial distribution of a
recycling source on a grid boundary (function STEP(ISTEP), see
section \ref{sec2.7.1}) as well as the distribution in velocity
space. The step function itself and the total source strength
FLUX(ISTRA), however, are defined in IF2COP, using data from input
block 14, see \ref{sec2.14}, for target recycling source ITARG. In
this case the sampling range from the spatial surface distribution
and the species distribution specified in input block 7 must be
equal to or a subrange of the the step-function for target
recycling source ITARG defined from the data in input block 14.

\medskip

if INDSRC(ISTRA) = 6, then the input flags described in section \ref{sec2.14}
are used to define the entire surface recycling source automatically. The
corresponding data in input block 7 for this stratum are not used and
may be omitted. Stratum ISTRA corresponds to the target recycling source
ITARG from block 14.

\subsection{entry IF3COP(ISTRAA,ISTRAE)}  \label{sec4.1.3}
{\sl Return data at the end of an EIRENE stratum}

\medskip

At this entry data are transferred back from EIRENE to the
interfacing module (and from there, after possibly further processing, to the external code). This entry is called from the ``strata-loop" in subr.
MCARLO, after all trajectories for a particular stratum ISTRA have
been sampled and after all volume and surface tallies have been
scaled and processed to their final form.

The call to IF3COP is controlled by the flag NMODE (input block 1).

At entry IF3COP to module INFCOP data are expected
for all strata ISTRA in the range ISTRA = ISTRAA,ISTRAE.
There they may be further prepared (e.g. normalized, or scaled to other units)
for transfer to the external code

\subsection{entry IF4COP} \label{sec4.1.4}
{\sl post-processing after one complete cycle: overall balances, statistics, etc.}

\section{Routines for cycling of EIRENE with external codes: EIRSRT}
\label{sec4.2}

In this subroutine the ``cycling" between EIRENE and external
codes is controlled. There are various options, such as
 ``full time dependent
(explicit)", ``quasi-stationary (explicit)", and
``quasi-stationary (implicit)".

\section{Routines for special tallies needed for code coupling:
UPTCOP}
\label{sec4.3}

Updating of code interface tallies COPV, ``along the fly". These tallies
are specific to the particular
background plasma CFD code, hence module UPTCOP is part of the code interfacing module.
Apart from scoring along Monte Carlo histories in this special routine
(rather than in EIRENE routine UPDATE),
the tally COPV is scaled, averaged, printed and plotted as any other EIRENE tally.

In older versions of EIRENE (before 2002) in particular momentum exchange rates (i.e. friction terms) in the
direction parallel to the magnetic field have been programmed here in UPTCOP.
Meanwhile these have become default tallies (i.e. scoring moved to UPDATE,
 see chapter \ref{sec.tables}. Still some particular
coupling tallies, e.g. to render the coupling procedure more implicit,
are scored in this routine.

\medskip

\newpage
\section{Statistical noise in Monte Carlo terms for external code,
noise-residuals: STATIS\_COP
}
\label{sec4.4}

As pointed out in Section \ref{sec2.9}:  for all primary (not derived)
EIRENE tallies the empirical standard deviation can also be obtained,
if requested, by setting proper flags in input block 9.


Furthermore: (since Nov. 2013): the new option UPFCOP, to obtain statistical estimates of linear
combination of standard tallies, by evaluating these composite scores after each trajectory
from the individual tally scores. This option is restricted to \textbf{linear} combinations of tallies, because
only then the sum over events and the sum over trajectories commute.

The original considerable CPU penalty for evaluating standard deviations along with
the mean values (the tallies themselves)
was avoided, at least for volume averaged tallies, in EIRENE
versions 99 and younger, by major code optimization (indirect addressing) in these
parts. In versions 2012 and younger the same optimization was carried out also for
standard deviations of all surface averaged tallies.

In Versions  Oct. 2013 and older some standard deviation tallies of special interest for code interfacing
(such as total, summed over species, source rates) have been implemented in a
special routine in the code interfacing module: STATIS\_COP.

This routine STATIS\_COP has now been made redundant, because all tallies needed for interfacing, including sums, differences of
default tallies, are now available as COPV tallies, either scored in UPTCOP or, for linear combinations,  in UPFCOP,
and for those COPV tallies standard deviations are available by default. This allows direct
and unbiassed estimation of
variances for source terms of internal (rather than total) energy balances, or ,
more generally, for plasma fluid equations
formulated in the rest frame of the plasma rather than in the laboratory frame.
For example the B2.5 version of the B2
2D plasma transport code contains the energy balance equations in the plasma frame,
but the momentum balance equations
in the laboratory frame. For all combinations direct statistical error estimates of
EIRENE sources are now available.


\bigskip

Version 2012 and older:


Subroutine \verb|STATIS_COP| provides these estimates for those
tallies which are specific to a particular coupled case. Usually
these will be source terms for particle, momentum, and energy
balance equations, summed over all donor species.

At entry \texttt{IF4COP} (see section \ref{sec4.1.4} above) these
may be further processed. For example, in case of coupling to the
B2 plasma fluid code these noise estimates are integrated into
global quantities (dimension: 1/time) and represent the
contribution of statistical noise to the overall residuals of a B2
run.

These are printed from \texttt{IF4COP}, together with the overall
balances of a coupled run (\texttt{TRCBAL=.true.})

Hence: if the B2 estimated residuals are of the same order as
these ``statistical noise residuals", then a further convergence
of the combined code can only be achieved by either increasing the
CPU-time for EIRENE.

Otherwise: The coupled B2-EIRENE run has failed to converge to the
solution within the statistical noise.

This is also represented by the convergence measure of ``saturated
residuals"
\medskip

\newpage
